\begin{afterword}
\summaryhead

The post names emergence as a present-day counterpart of phlogiston and vitalism. Emergence is usually defined, it says, as high-level behaviour arising from the interaction of many low-level elements, and taken literally that fits everything above the level of quarks. The author objects to the noun, not to the verb: ``X emerges from Y'' is fine when Y is a specific model with moving parts.

But ``emergence'' is commonly offered as an explanation in its own right, as in ``Intelligence is an emergent phenomenon!'', and in that use it meets every sign of a mysterious answer. It predicts nothing, has no moving parts, and is held with pride. Two exercises make the point. Deleting ``emergent'' from a sentence changes nothing, and replacing it with ``magical'' conveys the same knowledge. The word is popular, the author says, because it is ``the junk food of curiosity,'' and people with some science education use it as their ancestors used magic.

\responsehead

The post's narrowed point is fair: a word that labels a pattern is not a mechanism, and ``X emerges from Y'' says little until Y is a model. The problems are which sense of the word the post tests, and whom it holds responsible.

It tests the loosest definition available. It quotes the first sentence of a Wikipedia article whose next section gives narrower ones: Lewes's contrast between emergents and resultants, and the distinction between weak and strong emergence. In C. D. Broad's sense, a property is emergent only if it cannot, even in theory, be deduced from the properties of the parts. That is a claim that can turn out false, and the view lost ground as quantum chemistry and molecular biochemistry developed. The claim that the definition fits every phenomenon above quarks is true of the summary, not of the concept as those who used it defined it.

The two exercises, deleting the word or replacing it with ``magical,'' lose what the narrower senses of the word say. In C. D. Broad's usage, ``Life is an emergent phenomenon'' meant that life needs no special substance: Broad set emergent vitalism against the ``Entelechy'' of substantial vitalism. A life-substance is the error the previous post blamed vitalism for, and ``Life is a magical phenomenon'' does not rule it out. The post's own ``Even better'' account of the ant colony is close to Mark Bedau's definition of weak emergence. And the comparison with magic was already on the page the post links, where Bedau calls strong emergence ``uncomfortably like magic'' while defending the weak sense.

The people the post blames are not identified. The author has ``lost track'' of how many times the usage was heard, and quotes none of them; the charge that users ``take pride'' in their ignorance rests on nothing shown. Such writers have existed, Samuel Alexander among them, but the post names none. Its account of why the word is popular is also asserted without evidence.

In short: a fair point about a noun used in place of a mechanism, argued against the loosest definition on the page it cites, with two exercises that lose what the narrower senses of the word say, and with no user of the word quoted.
\end{afterword}
